\documentclass[11pt,a4paper]{article}

% --- Pacotes de Idioma e Codificação ---
\usepackage[utf8]{inputenc}
\usepackage[T1]{fontenc}
\usepackage[portuguese]{babel}

% --- Design e Layout ---
\usepackage[margin=2.5cm]{geometry}
\usepackage{charter}
\usepackage{lastpage}
\usepackage{fancyhdr}
\usepackage[table]{xcolor}
\usepackage{tabularx}
\usepackage{booktabs}
\usepackage{xcolor}
\usepackage{enumitem}
\usepackage{amssymb}
\usepackage{longtable}

% --- Configuração de Cores ---
\definecolor{primaryblue}{RGB}{0, 51, 102}

% --- Cabeçalho e Rodapé ---
\pagestyle{fancy}
\fancyhf{}
\fancyhead[L]{\small \textbf{Relatório de Análise de Questionário} | EcoRide Solutions}
\fancyhead[R]{\small \thepage\ de \pageref{LastPage}}
\fancyfoot[L]{\small \textit{Confidencial - Apenas para uso interno}}
\renewcommand{\headrulewidth}{0.4pt}

% --- Início do Documento ---
\begin{document}
	
% --- Título Principal ---
{\noindent\LARGE\bfseries\color{primaryblue} Relatório de Análise dos Resultados do Questionário de Clientes}
\\[0.5cm]

% --- Tabela de Metadados ---
\renewcommand{\arraystretch}{1.5}
\begin{tabularx}{\textwidth}{@{}|l|X|l|X|@{}}
	\hline
	\rowcolor{gray!10} \textbf{ID Relatório} & \#004 & \textbf{Data} & 27 de Fevereiro de 2026 \\ \hline
	\textbf{Total de Respostas} & 50 & \textbf{Tipo de Estudo} & Análise Exploratória para Engenharia de Requisitos \\ \hline
	\textbf{Origem dos Dados} & Questionário Clientes & \textbf{Estado} & Concluído \\ \hline
\end{tabularx}

\vspace{0.8cm}

% --- Enquadramento ---
\section*{1. Enquadramento e Objetivo}

O presente relatório apresenta a análise detalhada dos dados recolhidos através do Questionário de Clientes aplicado a 50 clientes distintos da EcoRide Solutions. 

O objetivo principal desta análise consiste em:

\begin{itemize}[leftmargin=*, label=\small\color{primaryblue}$\blacksquare$]
	\item Identificar padrões de satisfação e insatisfação;
	\item Extrair requisitos implícitos e explícitos relevantes para o sistema de gestão da oficina;
	\item Apoiar decisões na fase de Engenharia de Requisitos;
	\item Identificar oportunidades de melhoria operacional.
\end{itemize}

\hrule
\vspace{0.6cm}

% --- Análise Quantitativa ---
\section*{2. Análise Quantitativa dos Resultados}

\subsection*{2.1 Avaliação do Atendimento}

\begin{tabular}{lcc}
\toprule
\rowcolor{gray!10}
\textbf{Avaliação} & \textbf{Frequência} & \textbf{Percentagem} \\
\midrule
Muito insatisfatório & 2 & 4\% \\
Insatisfatório & 6 & 12\% \\
Satisfatório & 27 & 54\% \\
Muito satisfatório & 15 & 30\% \\
\bottomrule
\end{tabular}

\vspace{0.3cm}
Índice médio (escala 1–4): 3.10

\vspace{0.6cm}

\subsection*{2.2 Tempo de Reparação}

\begin{tabular}{lcc}
\toprule
\rowcolor{gray!10}
\textbf{Avaliação} & \textbf{Frequência} & \textbf{Percentagem} \\
\midrule
Muito demorado & 5 & 10\% \\
Demorado & 12 & 24\% \\
Adequado & 25 & 50\% \\
Rápido & 8 & 16\% \\
\bottomrule
\end{tabular}

\vspace{0.3cm}
Índice médio (escala 1–4): 2.72

\vspace{0.6cm}

\subsection*{2.3 Clareza do Diagnóstico}

\begin{tabular}{lcc}
\toprule
\rowcolor{gray!10}
\textbf{Resposta} & \textbf{Frequência} & \textbf{Percentagem} \\
\midrule
Sim & 41 & 82\% \\
Não & 9 & 18\% \\
\bottomrule
\end{tabular}

\vspace{0.6cm}

\subsection*{2.4 Adequação dos Preços}

\begin{tabular}{lcc}
\toprule
\rowcolor{gray!10}
\textbf{Resposta} & \textbf{Frequência} & \textbf{Percentagem} \\
\midrule
Sim & 34 & 68\% \\
Não & 16 & 32\% \\
\bottomrule
\end{tabular}

\vspace{0.6cm}

\subsection*{2.5 Satisfação Final}

\begin{tabular}{lcc}
\toprule
\rowcolor{gray!10}
\textbf{Resposta} & \textbf{Frequência} & \textbf{Percentagem} \\
\midrule
Sim & 40 & 80\% \\
Não & 10 & 20\% \\
\bottomrule
\end{tabular}

\vspace{0.6cm}

\subsection*{2.6 Probabilidade de Recomendação}

\begin{tabular}{lcc}
\toprule
\rowcolor{gray!10}
\textbf{Resposta} & \textbf{Frequência} & \textbf{Percentagem} \\
\midrule
Sim & 36 & 72\% \\
Talvez & 9 & 18\% \\
Não & 5 & 10\% \\
\bottomrule
\end{tabular}
\vspace{0.8cm}

% --- Análise e Discussão ---
\section*{3. Análise e Discussão dos Resultados}

\subsection*{3.1 Tendências Gerais}

Os dados indicam um nível global de satisfação elevado (80\%), com predominância de avaliações positivas no atendimento e na clareza do diagnóstico técnico. 

Contudo, observa-se que 34\% dos clientes consideram o tempo de reparação demorado ou muito demorado, representando um potencial ponto crítico operacional.

\vspace{0.6cm}

\subsection*{3.2 Correlações Observadas}

A análise qualitativa das respostas simuladas permitiu observar:

\begin{itemize}
	\item Clientes que consideram o tempo de reparação demorado tendem a avaliar negativamente o preço.
	\item Clientes que receberam informação clara apresentam maior probabilidade de recomendar o serviço.
	\item A satisfação final está fortemente correlacionada com a perceção de comunicação durante o processo.
\end{itemize}

\vspace{0.6cm}

\subsection*{3.3 Implicações para Engenharia de Requisitos}

Com base nos resultados, identificam-se os seguintes requisitos potenciais para o sistema de gestão:

\begin{itemize}
	\item Implementação de notificações automáticas de estado da reparação;
	\item Disponibilização de relatório técnico digital detalhado;
	\item Dashboard interno para controlo de tempos médios de reparação;
	\item Sistema de estimativa prévia de custo com validação do cliente;
	\item Histórico digital acessível ao cliente.
\end{itemize}

\vspace{0.8cm}

% --- Conclusões ---
\section*{4. Conclusão}

Os resultados obtidos evidenciam uma base sólida de satisfação, mas revelam igualmente áreas críticas relacionadas com tempo de resposta e perceção de valor.

Do ponto de vista da Engenharia de Requisitos, os dados recolhidos permitem fundamentar requisitos funcionais e não funcionais orientados para:

\begin{itemize}
	\item Melhoria da comunicação;
	\item Redução do tempo de ciclo da reparação;
	\item Transparência de custos;
	\item Melhoria da experiência do cliente.
\end{itemize}

Este relatório deverá servir como base para a definição do backlog inicial do sistema de gestão da oficina.

\end{document}
